\documentclass[10pt,a4paper]{amsart}
\usepackage[margin=19mm]{geometry}
\usepackage{amsmath,amssymb,mathtools}
\usepackage[scr=rsfs,cal=euler]{mathalfa}
\usepackage{xfrac}
\usepackage[unicode,hidelinks,bookmarks=false]{hyperref}
\newcommand{\ie}{i.e.\ }
\newcommand{\cf}{cf.\ }
\newcommand{\resp}{respectively\ }
\newcommand{\Subsection}{\subsection}
\providecommand{\nobreakdash}{\nobreak\hbox{-}}
\setlength{\emergencystretch}{2em}
\multlinegap0pt
\usepackage{amssymb}
\usepackage{bm}
\usepackage{enumerate}
\newtheorem{proposition}{Proposition}
\title{Manifold Data Imputation}
\author{David Levin}
\date{Source: arXiv:2604.12871v1 (CC BY 4.0)}
\begin{document}
\maketitle

\noindent\emph{Source and licence.} Manifold Data Imputation. Version arXiv:2604.12871v1 (CC BY 4.0), distributed by arXiv under the Creative Commons Attribution 4.0 International licence, http://creativecommons.org/licenses/by/4.0/, as recorded in the arXiv licence metadata for this exact version and shown on the abstract page https://arxiv.org/abs/2604.12871v1. Full licence notice: \texttt{licenses/arxiv-2604.12871-CC-BY-4.0.txt}.\par\smallskip

\noindent\textit{Editorial scope.} Counted unit: the article's proposition proposition (source0.tex lines 769--823). The article's complete original proof of it is delivered with it (source0.tex lines 825--875, one \texttt{proof} environment of 51 lines). No mathematics is written, repaired or re-derived: the statement and the proof are the article's own lines, in the article's own notation, and no retained line is substituted or re-flowed.  No further source-local context is needed: every cross-reference of every retained line resolves inside the two delivered spans.  Nothing was omitted from any retained span, so the package carries no omission record.  No retained line contains a citation, so the wrapper carries no bibliography and no bibliography entry.  No retained line contains a figure environment, an image-inclusion command or drawing code, so no image is delivered and the package declares no asset dependency.  Editorial formatting only, in addition: an AMS \texttt{amsart} preamble that restates the article's own declarations and macros used by the retained lines, namely the environments \texttt{proposition} on separate counters, together with the standard packages those lines need.  The retained environments print in this document as Proposition 1 in order of appearance, whereas the article prints the same items with the numbering of its own full document.  The complete native source of the article as distributed in the arXiv e-print for this exact version is bundled unchanged as \texttt{sources/198434/source0.tex}.
\begin{proposition}[Bounds on high-order differences
involving unknown data]
Let $u_h$ be the minimizer of
\[
\mathrm{J}(u)
=
\sum_{i\in G_k}\sum_{j=1}^d
\bigl(\Delta_j^{2k}u(i)\bigr)^2
\]
subject to $u_h|_X = f|_X$,
where $X$ is the set of known grid points,
and let $H$ denote the hole.

Assume that $f\in C^M(B)$,
with $M\ge 2k$.
Let $S_h$ be the set of all $(i,j)$
such that the stencil of $\Delta_j^{2k}u(i)$
contains at least one unknown point from $H$.

Then the following hold.

\begin{enumerate}
\item[(I)]
If the hole is contained in a ball
of diameter $qh$, then
\[
|S_h|\le C(q,d,k),
\]
and
\[
\bigl|\Delta_j^{2k}u_h(i)\bigr|
\le C\,h^{2k},
\qquad (i,j)\in S_h.
\]

\item[(II)]
If the hole is contained in a ball
of diameter $D$, independent of $h$,
then
\[
|S_h|\le C(D,d,k)\,h^{-d},
\]
and
\[
\bigl|\Delta_j^{2k}u_h(i)\bigr|
\le C\,h^{2k-d/2},
\qquad (i,j)\in S_h.
\]
\end{enumerate}

The constant $C$ depends on
$D$, $d$, $k$,
and $\|f\|_{C^{2k}(B)}$,
but is independent of $h$.
\end{proposition}

\begin{proof}
Since $u_h$ minimizes $\mathrm{J}$,
and the exact data $f$ define an admissible extension,
we have
\[
\sum_{i,j}
\bigl(\Delta_j^{2k}u_h(i)\bigr)^2
\le
\sum_{i,j}
\bigl(\Delta_j^{2k}f(i)\bigr)^2.
\]
In particular,
\[
\sum_{(i,j)\in S_h}
\bigl(\Delta_j^{2k}u_h(i)\bigr)^2
\le
\sum_{(i,j)\in S_h}
\bigl(\Delta_j^{2k}f(i)\bigr)^2.
\]

Since $f\in C^{2k}(B)$,
\[
\bigl|\Delta_j^{2k}f(i)\bigr|
\le C\,h^{2k},
\]
and hence
\[
\bigl(\Delta_j^{2k}f(i)\bigr)^2
\le C\,h^{4k}.
\]

In case (I),
the number of affected stencils is bounded,
so
\[
\sum_{(i,j)\in S_h}
\bigl(\Delta_j^{2k}f(i)\bigr)^2
\le C\,h^{4k}.
\]
In case (II),
the number of affected stencils is $O(h^{-d})$,
so
\[
\sum_{(i,j)\in S_h}
\bigl(\Delta_j^{2k}f(i)\bigr)^2
\le C\,h^{4k-d}.
\]

Since each term is bounded by the total sum,
the stated pointwise bounds follow.
\end{proof}

\end{document}
